% Stage 8 controlled rewrite. No new data, estimator, hypothesis, or result.

\begin{abstract}
\noindent We examine whether the informational association between digital-payment
activity and formal borrowing, documented at the firm level, generalises to individuals.
Using Global Findex repeated cross-sections for a common frame of 93 economies observed in
2021 and 2024 --- a sample and specification developed during exploratory analysis of the
data rather than fixed by a prospective pre-registration --- we estimate survey-weighted
linear probability models with economy fixed effects and economy-clustered inference.
Digital-payment activity is positively associated with formal borrowing in both waves, and
its interaction with lower 2019 historical credit-reporting coverage is negative in 2021
($-1.74$ percentage points; wild-cluster $p=0.002$) and 2024 ($-1.61$ percentage points;
wild-cluster $p=0.020$) --- the opposite of the sign a credit-information-substitution
mechanism predicts. We then subject this reversal to two checks a genuine substitution
channel should survive, and it survives neither. An identical interaction estimated on
informal borrowing, a channel with no formal credit-information screening, is nonetheless
positive and significant in both waves ($+1.95$, $+1.75$ percentage points) --- evidence
against a formal-credit-specific, information-screening interpretation of the reversal,
though informal borrowing is not a clean negative control (it can still reflect social
ties, reputation and selection). Allowing digital-payment activity to also interact with
log GDP per capita and the account-ownership rate erodes the formal-borrowing interaction
to non-significance in both waves, though this is consistent with severe collinearity
between coverage and GDP per capita ($r=0.74$ across the 93 economies) rather than with a
cleanly identified competing channel. The result is an empirical puzzle rather than a
confirmed boundary condition: a cross-level association observed consistently across both
waves, whose two most natural explanations --- an information-screening channel specific
to formal credit, and a general economic-development channel --- are, respectively,
disfavoured by the evidence and not separately identifiable with this design. An
inverse-probability-reweighting check further shows the 93-economy frame cannot be
extrapolated to the higher-coverage economies the Findex module excluded. We identify no
mechanism; account access, lender processing capacity, supply-side rationing, measurement
and firm--retail lending-technology differences remain open explanations.

\medskip
\noindent\textbf{Keywords:} digital payments, formal borrowing, credit information,
financial inclusion, Global Findex, external validity, placebo test, exploratory
analysis.
\quad\textbf{JEL:} G21, G51, O16, D14.
\end{abstract}

\section{Introduction}

Digital payments can generate records that lenders may use when assessing borrowers. The
economic question is not only whether payment activity correlates with credit access, but
whether the informational value of that activity travels across lending contexts. The
closest motivating benchmark is the firm-level study of Galilea, Farazi and Mare, which
reports weaker firm credit constraints among electronic-payment users and a larger
association in economies with weaker credit-information systems
\citep{GalileaFaraziMare2026}. Their result motivates a substitution prediction: when
conventional information is scarce, payment records should have greater marginal value.

Firms and individuals, however, do not face the same lenders, products, information
technology or access margins. A firm outcome such as a self-reported credit constraint is
also different from an individual's realised formal borrowing. We therefore test whether
the conditional pattern generalises to individuals using Global Findex evidence, on a
common-economy specification developed during exploratory analysis of the data rather
than fixed in advance of estimation. The 2021 and 2024 waves are repeated cross-sections, not an individual panel;
2017 is retained only as contextual/supplementary evidence because its payment measure is
not a strict common exposure.

The primary result is a contrary one. Digital-payment activity is positively associated
with formal borrowing, but the association is weaker in economies with thinner historical
2019 formal credit-reporting coverage. The inverse gradient is observed in both 2021 and
2024 under a qualified-comparable payment measure. The result is not confined to one
economy under the leave-one-economy-out check and survives wild-cluster inference. This
reversal is then run through two diagnostic checks a genuine formal-credit-specific
information channel should survive --- an informal-borrowing placebo and a set of
macro-control interactions, both described in \S5.4 --- and it survives neither cleanly.

Our contribution is therefore not a boundary condition in the usual sense but evidence
against a specific mechanism. The account motivating the search for this reversal ---
that digital-payment records substitute for scarce formal credit information in a way
specific to formal, information-screened lending --- does not survive a placebo test on
informal borrowing, a channel with no such screening, though informal borrowing is not a
clean negative control (it can still reflect social ties, reputation and selection). The
account's status once general economic development is accounted for cannot be separated
from a design limitation rather than confirmed as a competing channel. We do not
claim that the firm-level result is false, nor do we identify what does explain the
individual-level reversal. The evidence remains compatible with account-access selection,
general economic development, lender processing capacity, supply-side rationing, record
usability, measurement and differences between firm and retail lending technologies ---
an empirical puzzle rather than a resolved mechanism.

\section{Related literature and positioning}

\paragraph{Global financial-inclusion measurement.} The individual-level data underlying
this paper come from the Global Findex Database, the household-survey program that
underpins most cross-country evidence on account ownership, digital payments and
borrowing behaviour
\citep{DemirgucKuntKlapperSingerAnsar2022,KlapperSingerStaritaNorris2025}. Those reports
document the digital-payment and formal-borrowing series used here and the module and
country-coverage rules that generate the 93-economy common frame (\S3); they do not
themselves test a credit-information-substitution mechanism.

\paragraph{Firm and SME finance.} Galilea, Farazi and Mare study firm credit constraints
and electronic payments across economies, using Business Ready credit-information
dimensions as institutional moderators \citep{GalileaFaraziMare2026}. Their unit, outcome,
moderator and design differ from ours. The present paper is consequently an
external-population and cross-level comparison, not a direct replication.

Experimental firm evidence also shows that lowering payment-adoption barriers can increase
access to mobile loans for SMEs \citep{DaltonPamukRamrattanUrasVanSoest2024}. This supports
the economic relevance of payment technology in business finance, but does not establish
the same boundary for retail borrowing.

\paragraph{Digital footprints and lender information.} Digital footprints can complement
rather than substitute for bureau information in predicting consumer default in a
single-lender setting \citep{BergBurgGombovicPuri2020}. Recent lender-level evidence uses
applications, payment histories, approval, pricing and default outcomes to study the
informational content of cashless payments \citep{Ghosh2026}. These studies clarify what
is needed to identify information value or lender use; the Findex design does not observe
those links.

\paragraph{Positioning.} We make no categorical first claim. The paper contributes a
transparent individual-level test, developed through exploratory analysis of the data and
subsequently independently reproduced (Stage 12), whose inverse gradient is observed in
both recent waves. It is a descriptive boundary-condition result and a theory-generating
puzzle, not an identified mechanism or causal institutional effect.

\section{Data and measurement}

The primary frame contains 93 economies common to the 2021 and 2024 Global Findex waves
\citep{WorldBankFindexMicrodata2025,DemirgucKuntKlapperSingerAnsar2022,KlapperSingerStaritaNorris2025}
and 194,558 respondents. The 2024 baseline is 97, not the raw release's 140, because the
borrowing/digital-payment module was administered in only 97 economies; the 93-economy
frame is the intersection of that 97-economy set with the 139 economies in the 2021 raw
release, after also requiring complete primary covariates. This module-administration rule
is not neutral with respect to the moderator: of the 133 economies present in both raw 2021
and 2024 releases, the 40 excluded from the 93-economy frame are overwhelmingly high-income
economies (for example Australia, Canada, France, Germany, Japan, Korea, Switzerland, the
United Kingdom and the United States) with near-universal 2019 credit-information coverage
--- 24 of the 34 economies worldwide with 2019 coverage at or above 95\% are excluded from
the 93-economy frame, and the moderator's interquartile range is compressed in-sample (11.7
to 60.7) relative to the excluded-but-data-available pool (1.2 to 92.6). The 93-economy
sample is therefore not a random or exhaustive cross-section of economies with available
2019 coverage data; it is a population pre-filtered by Findex fieldwork decisions that are
plausibly correlated with the same financial-system completeness the moderator is meant to
capture. Findex waves are repeated cross-sections. The 2021 release's interview-year field
includes some observations interviewed in 2022; they remain part of the official 2021
survey wave and are not silently reclassified.

The outcome (\texttt{fin22a}) is a single self-reported item --- ``In the past 12 months,
have you borrowed any money from a bank or a similar financial institution?'' --- rather
than a constructed aggregate of named credit products; its ``or a similar financial
institution'' phrasing leaves some interpretive latitude to respondents about what counts
as formal, which this paper does not attempt to resolve further. The exposure is the
qualified-comparable 2021--2024 digital-payment measure documented in the harmonisation
audit. It is not described as an invariant causal treatment. The 2017 payment
reconstruction is contextual only.

The moderator is the maximum of private-bureau and public-registry coverage in 2019,
standardised so that larger values denote thinner historical formal credit-reporting
coverage. It measures breadth of recorded coverage, not information depth, accessibility,
legal usability or actual lender use. Coverage, depth and accessibility therefore remain
separate constructs. This series is drawn from the archived World Bank Doing Business
``Getting Credit'' indicators \citep{WorldBankDoingBusinessArchive} (the same 2019-vintage
snapshot used in the frozen \texttt{H-000500}/\texttt{H-000501} single-wave design), which
the World Bank paused and later discontinued in 2021 after an internal review found data
irregularities in a subset of country scores \citep{WorldBank2020Irregularities}, publicly
confirmed for China (DB2018) and Saudi Arabia (DB2020) ---
both present in this 93-economy sample. Whether the archived 2019 values used here reflect
pre- or post-review figures for the affected economies is not established from the files
available to this project; this is disclosed as a provenance uncertainty, not a
demonstrated error, consistent with the same disclosure already carried in the
\texttt{H-000500}/\texttt{H-000501} manuscript. Excluding China and Saudi Arabia and
re-standardizing the moderator over the remaining 91 economies leaves the primary
interaction negative and significant in both waves ($-1.64$ pp in 2021, wild-cluster
$p=0.001$; $-1.54$ pp in 2024, wild-cluster $p=0.007$; changes of 5.7\% and 4.3\%
relative to the 93-economy estimates), so the registered reversal is not driven by the two
in-sample economies with documented \emph{Getting Credit} irregularities. This does not
validate the Doing Business production process for the remaining 91 economies or rule out
general moderator measurement error.

Account ownership is retained as an access/selection-sensitive variable. The primary
specification does not condition on it. Conditioning or restricting to account holders is
a sensitivity estimand that may remove part of the access pathway or create collider risk;
it is not a universally superior control specification.

\section{Empirical strategy}

For each wave $w\in\{2021,2024\}$, the primary model is a survey-weighted linear
probability model with economy fixed effects:
\[
Y_{icw}=\beta_{1w}D_{icw}+\beta_{3w}(D_{icw}\times L_c^{2019})
+X_{icw}\gamma_w+\delta_{cw}+\varepsilon_{icw},
\]
where $D$ is qualified-comparable digital-payment activity, $L^{2019}$ is the fixed
historical low-coverage moderator, and $X$ contains the frozen demographic controls: a
female indicator, a centered age quadratic (age and age$^2$), and categorical education
and income-quintile indicators.
Weights are the official wave-specific weights normalised within economy-wave; no pooled
survey weight is created. Inference clusters at the economy level. Interaction inference
uses a null-imposed, cluster-score Rademacher bootstrap (999 replications) that
reconstructs the coefficient distribution from a fixed bread matrix computed once at the
unrestricted fit, rather than refitting the model on every replicate; this is referred to
below as the wild-cluster bootstrap, following common usage, though it is a one-step
score-bootstrap variant and not a full-refit Cameron--Gelbach--Miller wild-cluster
bootstrap.

The pooled model is secondary and descriptive. It includes wave indicators and the
payment-by-wave and payment-by-low-coverage-by-wave terms. Since $L_c^{2019}$ does not
change between 2021 and 2024, the triple interaction is a cross-wave conditional contrast,
not an effect of institutional change.

The primary estimands are the wave-specific association and wave-specific coverage
gradient. The difference between gradients is described only as a temporal contrast;
coefficient equality is not asserted because no new equality test was pre-specified.

\section{Results}

\subsection{Primary two-wave evidence}

Digital-payment activity is positively associated with formal borrowing in both waves.
The payment-by-low-coverage interaction is negative in both primary separate-wave models:
$-1.74$ percentage points in 2021 and $-1.61$ percentage points in 2024. The wild-cluster
bootstrap p-values are 0.002 and 0.020, respectively. These are conditional associations,
not causal effects. Table~\ref{tab:stage8primary} reports the primary estimates.

\begin{table}[htbp]
\centering
\caption{Primary separate-wave estimates, 2021 and 2024}
\label{tab:stage8primary}
\begin{adjustbox}{max width=\textwidth}
\begin{tabular}{lrr}
\toprule
 & 2021 & 2024 \\
\midrule
Digital payment association (pp) & 9.07*** & 8.13*** \\
Payment $\times$ lower 2019 coverage (pp) & -1.74*** & -1.61*** \\
Clustered SE, interaction (pp) & (0.49) & (0.60) \\
Wild-cluster $p$, interaction & 0.002 & 0.020 \\
Observations & 96,711 & 97,847 \\
Economies & 93 & 93 \\
\bottomrule
\end{tabular}
\end{adjustbox}
\smallskip
\begin{minipage}{\linewidth}\footnotesize \emph{Notes.}
Survey-weighted WLS linear probability models with economy fixed effects and
economy-clustered standard errors. The exposure is qualified-comparable across
2021--2024, not asserted to be perfectly invariant. Coefficients are percentage
points. Wild-cluster inference uses 999 null-imposed cluster-score Rademacher
replications. *** denotes clustered $p<0.01$, ** $p<0.05$, * $p<0.10$.
Generated by \texttt{code/\allowbreak{}study\_05\_crosscountry/\allowbreak{}tables/\allowbreak{}build\_stage8\_tables.py}
from \texttt{results/\allowbreak{}stage5\_estimation/\allowbreak{}estimates.csv} and
\texttt{results/\allowbreak{}stage5\_1\_reconciliation/\allowbreak{}wild\_cluster\_bootstrap.csv}.\end{minipage}
\end{table}


The repeated appearance of the inverse gradient makes a one-year anomaly explanation less
plausible. But two-wave persistence does not distinguish a recurring economic relationship
from a recurring selection structure, measurement structure or institutional composition.
It does not prove temporal stability, equality, a structural effect or a common mechanism.
Every one of the 93 leave-one-economy-out estimates remains negative in each
wave, with ranges of approximately $[-1.97,-1.56]$ pp in 2021 and $[-1.83,-1.27]$ pp in
2024. Because the primary weighting gives each economy equal aggregate weight, dropping
any single economy removes only about one ninety-third of total estimation weight, which
mechanically narrows this range; sign preservation across all 93 single-economy deletions
rules out a dominant-influence economy, not homogeneity of the underlying association
across economies or economy groups.

\subsection{Pooled contrast and reconciliation}

The pooled secondary interaction is $-1.98$ percentage points. The 2024 triple interaction
is $+0.58$ percentage points, with wild-cluster $p=0.305$ (Table~\ref{tab:stage8pooled}).
Thus the specified pooled model finds no statistically detectable wave difference in the
conditional gradient. This is not proof that the two coefficients are identical.

\begin{table}[htbp]
\centering
\caption{Secondary pooled repeated-cross-section parameterization}
\label{tab:stage8pooled}
\begin{tabular}{lr}
\toprule
Term & Estimate (pp) \\
\midrule
Payment $\times$ lower 2019 coverage & -1.98*** \\
Payment $\times$ lower coverage $\times$ 2024 & +0.58 \\
Wild-cluster $p$, triple interaction & 0.305 \\
Observations & 194,558 \\
Economies & 93 \\
\bottomrule
\end{tabular}
\smallskip
\begin{minipage}{0.92\linewidth}\footnotesize \emph{Notes.}
This is a descriptive pooled parameterization. The 2019 moderator is fixed across
waves, so the triple interaction is a temporal contrast and not an institutional-change
effect. The pooled-implied gradients need not equal the separate-wave coefficients
because nuisance slopes and interaction terms are jointly estimated on the pooled sample.
*** denotes clustered $p<0.01$, ** $p<0.05$, * $p<0.10$; the triple interaction is not
significant at conventional levels. Generated by
\texttt{code/\allowbreak{}study\_05\_crosscountry/\allowbreak{}tables/\allowbreak{}build\_stage8\_tables.py} from
\texttt{results/\allowbreak{}stage5\_estimation/\allowbreak{}estimates.csv} and
\texttt{results/\allowbreak{}stage5\_1\_reconciliation/\allowbreak{}wild\_cluster\_bootstrap.csv}.
\end{minipage}
\end{table}


The pooled-implied 2021 gradient is $-1.98$ pp and the pooled-implied 2024 gradient is
$-1.40$ pp. They do not equal the separate-wave estimates mechanically because the pooled
regression jointly estimates the interaction terms with common nuisance slopes, pooled
sample composition and the pooled fixed-effects parameterisation. The separate-wave
estimates remain primary.

\subsection{Institutional dimensions and account sensitivity}

Depth is analysed separately from coverage and is secondary evidence. Accessibility is a
validation/context layer rather than a comparable primary moderator in the frozen dataset.
The evidence therefore does not establish a general institutional gradient across all
dimensions.

The account-ownership sensitivity is read only as selection/access sensitivity. Conditioning
on account ownership slightly attenuates the interaction in 2021 and makes it more negative
in 2024. This cross-wave difference does not identify account access as the mechanism. In
the 2021 wave specifically, approximately 3\% of \texttt{account\_fin}\,$=1$ respondents are
constructed from payment- or card-derived evidence rather than the standard
account-screening questions, so \texttt{account\_fin} and \texttt{digital\_payment} share a
small construction overlap in that wave --- a first-order collinearity risk for the
account-conditioning specification in 2021 specifically, distinct from the
selection-versus-mechanism caveat above.

\subsection{Additional robustness: external-validity reweighting and macro-control interactions}

Two further exploratory checks quantify, rather than only narrate, two threats already
raised above and in the Limitations: the module-administration sample restriction, and the
possibility that the coverage moderator proxies general economic or financial development.
Neither check is pre-registered or confirmatory; both are diagnostic, following the same
inverse-probability-weighting protocol already used for the frozen H-000501 record's Q5
module-selection check, adapted to the present two-wave frame.
Table~\ref{tab:stage8robustness} reports both.

\emph{External-validity reweighting.} Restricting to the 132 economies with usable
economy-level covariates among the 133 present in both raw 2021 and 2024 releases (Taiwan
is dropped for missing all covariates, as in the frozen H-000501 precedent), an inclusion
logit on 2019 coverage, log GDP per capita, region and log adult population separates the
93 included from the 39 excluded economies with a $c$-statistic of 0.982 (5-fold
cross-validated mean AUC 0.954), indicating the near-complete separation reflects a
genuinely strong inclusion rule --- not overfitting on 132 economies with four
predictors --- so it is read as a description of the selection process rather than a
numerically unstable fit; inverse-probability weights are trimmed at the 1st/99th
percentiles (maximum trimmed weight 6.44, Kish effective sample size 60.9 of 93) before
reweighting. Reweighting the
93-economy sample by the inverse inclusion probability keeps the coverage interaction
negative and of similar magnitude in 2021 (primary $-1.74$ pp; reweighted $-1.77$ pp, a
1.3\% change, wild-cluster $p=0.002$) and attenuates it further in 2024 (primary $-1.61$
pp; reweighted $-1.23$ pp, a 23.3\% change, wild-cluster $p=0.043$), in both cases well
inside the $\pm 50\%$ stability band. The classification is nonetheless
\textbf{not credibly transportable with the available covariate overlap}, not because the
reweighted point estimates are unstable, but
because 80 of the 132 economies fall outside common support --- the excluded, mostly
high-coverage economies are sufficiently different in covariate profile from the included
93 that no amount of reweighting can credibly extrapolate to them. This is the same
qualitative outcome, and a similar outside-support count, as the frozen H-000501 Q5 check
on a different target frame, which strengthens the reading that the module-administration
restriction is a genuine, not incidental, limit on external validity.

\emph{Macro-control interactions.} Because \texttt{lowcov\_z} interacted with log GDP per
capita or the account-ownership rate is collinear with the economy fixed effects already in
M2 (both terms are constant within economy), the diagnostic instead interacts
\texttt{digital\_payment} with these two economy-level covariates --- the same construction
already used for the primary interaction --- and asks whether \texttt{dig\_x\_lowcov}
survives. It does not survive at conventional significance: the coefficient attenuates from
$-1.74$ to $-1.21$ pp in 2021 (30.9\% change, $p=0.087$, no longer significant at 5\%) and
from $-1.61$ to $-0.98$ pp in 2024 (39.0\% change, $p=0.31$, not significant), while
remaining negative in sign in both waves. Neither new interaction term itself reaches
conventional significance (payment $\times$ log GDP per capita: $p=0.75$ and $p=0.64$;
payment $\times$ account-ownership rate: $p=0.14$ and $p=0.07$). Part of this attenuation
is mechanical rather than substantive: the 2019 coverage moderator and log GDP per capita
are themselves highly correlated across the 93 economies ($r=0.74$, $R^2=0.55$), so asking
the model to identify both the coverage and the GDP-per-capita interaction terms
simultaneously, with only 93 economy-level clusters, imposes a severe collinearity burden
regardless of whether a genuinely separable development channel exists --- consistent with
neither new interaction term itself reaching significance. Excluding China and Saudi
Arabia and re-standardizing the coverage and GDP-per-capita moderators over the remaining
91 economies does not change this reading: the coverage interaction is $-1.38$ pp
($p=0.039$) in 2021 and $-0.95$ pp ($p=0.33$) in 2024, and \texttt{dig\_x\_lgdppc} and
\texttt{dig\_x\_accrate} remain non-significant in both waves, so the attenuation pattern
is not driven by the two economies with disclosed \emph{Getting Credit} irregularities. This is reported plainly rather
than softened: with this degree of correlation and this many clusters, the design cannot
cleanly distinguish a credit-information-specific channel from a general economic-development
channel. The attenuation is consistent with either genuine confounding or with
collinearity-driven loss of precision, and the data cannot adjudicate between the two. We
therefore do not claim the coverage interaction is spurious, only that its specificity to
credit-information coverage, as opposed to economic development more broadly, is not
identified by this design.

\begin{table}[htbp]
\centering
\caption{Stage 17 additional robustness: external-validity reweighting and macro-control interactions}
\label{tab:stage8robustness}
\begin{adjustbox}{max width=\textwidth}
\begin{tabular}{lrr}
\toprule
\multicolumn{3}{l}{\textit{Panel A: IPW reweighting toward the 39 excluded, mostly high-coverage economies}} \\
 & 2021 & 2024 \\
\midrule
Primary $\times$ lower-coverage interaction (pp) & -1.74 & -1.61 \\
IPW-reweighted interaction (pp) & -1.77 & -1.23 \\
Clustered SE, reweighted (pp) & (0.51) & (0.64) \\
Wild-cluster $p$, reweighted & 0.002 & 0.043 \\
Relative change vs.\ primary & 1.3\% & 23.3\% \\
\addlinespace
\multicolumn{3}{l}{Inclusion logit $c$-statistic: 0.982; economies outside common support: 80 of 132} \\
\multicolumn{3}{l}{Max IPW weight after 1st/99th trim: 6.44; Kish ESS: 60.9 of 93} \\
\multicolumn{3}{l}{\textbf{Classification: not transportable}} \\
\addlinespace
\multicolumn{3}{l}{\textit{Panel B: digital-payment interactions with log GDP per capita and account-ownership rate}} \\
 & 2021 & 2024 \\
\midrule
Payment $\times$ lower coverage (pp), with macro controls & -1.21* & -0.98 \\
 & (0.70) & (0.97) \\
Payment $\times$ log GDP per capita (pp) & -0.24 & -0.35 \\
 & (0.76) & (0.74) \\
Payment $\times$ account-ownership rate (pp) & 1.17 & 1.29* \\
 & (0.79) & (0.73) \\
\addlinespace
\multicolumn{3}{l}{Primary (no macro controls) interaction (pp): -1.74 (2021), -1.61 (2024)} \\
\bottomrule
\end{tabular}
\end{adjustbox}
\smallskip
\begin{minipage}{\linewidth}\footnotesize \emph{Notes.}
Panel A: inverse-selection weighting ($1/\hat p$, trimmed at the 1st/99th percentile,
combined multiplicatively with the primary economy-equal weight), following the
same protocol as \texttt{H-000501}'s Q5 module-selection check, adapted to a
different target frame (economies present in both raw 2021 and 2024 Findex releases).
Classification follows the same three-way rule: \emph{stable} requires adequate common
support; the 80-of-132 economies outside common support here
force \emph{not transportable} regardless of the reweighted point estimate's own stability.
Panel B: digital\_payment $\times$ log-GDP-per-capita and digital\_payment $\times$
account-ownership-rate terms are added \emph{in place of} lower-coverage interactions
with these macro covariates, because the latter are collinear with the economy fixed
effects already in M2 (both a covariate constant within economy and its interaction with
another economy-constant covariate are absorbed by \texttt{C(economycode)}). *** denotes
clustered $p<0.01$, ** $p<0.05$, * $p<0.10$. Both panels are exploratory/diagnostic
robustness, not pre-registered or confirmatory. Generated by
\texttt{code/\allowbreak{}study\_05\_crosscountry/\allowbreak{}tables/\allowbreak{}build\_stage8\_robustness\_table.py} from
\texttt{results/\allowbreak{}stage17\_external\_validity/\allowbreak{}}.
\end{minipage}
\end{table}


\subsection{Placebo check: informal borrowing}

A further exploratory check asks whether the coverage interaction is specific to formal
borrowing or also appears for an outcome the substitution reading predicts should be
unaffected: borrowing from family or friends (\texttt{informal\_borrow}, from
\texttt{fin22b}), constructed exactly as in the frozen H-000500 record and re-estimated
under the identical M2 specification, weighting and wild-cluster inference used
throughout. The row-level sample reconstruction was validated against
\texttt{results/stage5\_estimation/estimates.csv} before this check was read: both waves'
complete-case $N$ and $G$ match the primary formal-borrowing sample exactly
(96,711/93 in 2021, 97,847/93 in 2024), confirming the reconstruction is faithful;
\texttt{informal\_borrow}'s own missingness (dropping \texttt{fin22b} DK/refused) then
yields a slightly smaller estimation sample (96,535 and 97,723 respondents, still 93
economies in both waves). Table~\ref{tab:stage8placebo} reports both waves.

The placebo does not pass. As with the frozen H-000500 record's own informal-borrowing
check, the coverage interaction on informal borrowing is \textbf{positive}, not
approximately zero: $+1.95$ percentage points in 2021 (wild-cluster $p=0.009$) and $+1.75$
percentage points in 2024 (wild-cluster $p=0.006$) --- opposite in sign to the
$-1.74$/$-1.61$ percentage-point formal-borrowing interaction, and, unlike H-000500's
single-wave version of this check, now replicated with a similar magnitude and precision
in both waves of a fully separate design. This is reported plainly rather than
reconciled: the data are consistent with a reallocation of borrowing activity toward
informal sources where coverage is thinner, but this check alone does not establish that
reading, since it does not observe the same individuals borrowing from both sources or
rule out a shared measurement or selection process behind both outcomes. It does,
however, further undermine any reading of the formal-borrowing interaction as confirming
the credit-information-substitution prediction: a genuine substitution effect specific to
formal, information-screened lending should not also show up, with similar strength and
in the same two waves, for borrowing that involves no credit-information screening at
all. Excluding China and Saudi Arabia and re-standardizing the moderator over the
remaining 91 economies leaves this placebo interaction positive and significant in both
waves ($+1.81$ pp in 2021, $+1.69$ pp in 2024, both wild-cluster $p<0.03$; changes of 7.1\%
and 3.9\% relative to the 93-economy estimates), so it is not an artifact of the two
economies with disclosed \emph{Getting Credit} irregularities either.

\begin{table}[htbp]
\centering
\caption{Placebo check: informal borrowing, 2021 and 2024}
\label{tab:stage8placebo}
\begin{adjustbox}{max width=\textwidth}
\begin{tabular}{lrr}
\toprule
 & 2021 & 2024 \\
\midrule
Digital payment association (pp) & +9.71*** & +9.50*** \\
Payment $\times$ lower 2019 coverage (pp) & +1.95*** & +1.75*** \\
Clustered SE, interaction (pp) & (0.72) & (0.59) \\
Wild-cluster $p$, interaction & 0.009 & 0.006 \\
Observations & 96,535 & 97,723 \\
Economies & 93 & 93 \\
\bottomrule
\end{tabular}
\end{adjustbox}
\smallskip
\begin{minipage}{\linewidth}\footnotesize \emph{Notes.}
Outcome is \texttt{informal\_borrow} (borrowed from family or friends in the past 12
months), constructed from \texttt{fin22b} exactly as in the frozen H-000500 record.
Same M2 specification, weighting and null-imposed cluster-score Rademacher wild-cluster
inference (999 replications) as the primary formal-borrowing model. *** denotes clustered
$p<0.01$, ** $p<0.05$, * $p<0.10$. Generated by
\texttt{code/\allowbreak{}study\_05\_crosscountry/\allowbreak{}tables/\allowbreak{}build\_stage8\_placebo\_table.py} from
\texttt{results/\allowbreak{}stage18\_placebo/\allowbreak{}informal\_borrow\_estimates.csv}.\end{minipage}
\end{table}


\section{Discussion}

The central empirical fact is that the inverse digital-payment/formal-borrowing gradient is
observed in both recent Findex waves. The two-wave result reduces the likelihood of a
one-wave anomaly, but it does not by itself identify an underlying economic mechanism: the
same persistent selection or measurement structure could recur.

Two-wave persistence, however, is not the paper's strongest result. Two checks were run
specifically to ask whether the reversal reflects the credit-information-substitution
mechanism the literature motivates, and both point away from it. The informal-borrowing
placebo (\S5.4) finds the identical interaction, on a channel with no formal
credit-information screening, positive and significant in both waves --- the opposite of
what a mechanism specific to formal, information-screened lending would produce. This is
evidence against that specific interpretation, not merely a failure to confirm it: the
formal-credit-specific account predicts a near-zero placebo interaction, and instead the
placebo interaction is of similar magnitude and precision to the formal-borrowing result,
with the opposite sign --- though informal borrowing is not a clean negative control, since
it can still reflect social ties, reputation and selection unrelated to credit-information
screening. Separately, allowing digital-payment
activity to interact with log GDP per capita and the account-ownership rate (\S5.4) erodes
the formal-borrowing interaction to non-significance in both waves; because 2019 coverage
and log GDP per capita are themselves correlated at $r=0.74$ across the 93 economies, this
is consistent with either a genuine development-level channel or with the two variables
simply being too collinear, given 93 clusters, for the design to tell them apart. Taken
together, these checks leave the reversal itself observed consistently across both waves,
but its two most natural explanations --- a credit-information-specific channel, and a
general-development channel --- respectively disfavoured by the evidence and unidentifiable.
The paper's contribution is therefore best read as an empirical puzzle: a pattern that
recurs across both waves and resists the mechanism motivating the search for it, not a
confirmed boundary condition with an unresolved but plausible cause.

This finding also describes a narrower population than ``cross-country'' evidence might
suggest. The 93-economy frame excludes most near-universal-coverage economies because the
underlying survey module was not administered there, not because moderator data are
missing (see Data); the boundary-condition and transportability language in this paper
should therefore be read as describing module-administered economies with predominantly
low-to-middle 2019 credit-information coverage, not economies in general.

Several explanations remain open. First, in thin formal systems, account ownership and
digital-payment use may mark entry into formal finance rather than provide an incremental
lender-legible signal. Second, retail lenders may lack the systems to process payment
records even when such records exist. Third, realised borrowing combines borrower demand,
selection and lender supply; a survey outcome cannot separate them. Fourth, coverage is an
imperfect historical proxy and payment measures are qualified-comparable rather than
identical. Finally, firm and retail lending technologies may use payment information in
different ways.

The present evidence does not choose among the remaining explanations, but it does narrow
them. A channel operating specifically through formal-lender use of payment-based credit
information is difficult to reconcile with a placebo interaction of similar size and
significance on a channel where no such information use is possible; that specific
account is therefore not merely unconfirmed but actively disfavoured by the evidence.
Account-access selection, general economic development, lender processing capacity,
supply-side rationing, record usability, measurement and firm--retail lending-technology
differences all remain compatible with what is observed. Direct mechanism evidence would
require applications, approval/rejection, lender type, loan terms, account-opening timing,
payment-record access or actual underwriting use.

The result also does not overturn the firm-level literature. The benchmark differs in unit,
outcome, moderator and lending technology. The appropriate inference is narrower: a
firm-level substitution pattern should not be assumed to generalise automatically to
individual formal borrowing.

\section{Conclusion}

Digital-payment activity is positively associated with formal borrowing in the 2021 and
2024 Global Findex repeated cross-sections, and its interaction with thinner historical
2019 formal-credit-reporting coverage is negative in both waves --- the opposite of what
the motivating credit-information-substitution prediction implies.

The paper's contribution is evidence against a specific mechanism rather than a confirmed
boundary condition. An informal-borrowing placebo that a formal-credit-specific
information-screening channel should have left near zero instead returns a positive,
significant interaction of similar magnitude in both waves --- evidence against that
specific account, though informal borrowing is not a clean negative control, since it can
still reflect social ties, reputation and selection. A parallel check allowing
digital-payment activity to interact with general economic development (log GDP per
capita) and the account-ownership rate erodes the formal-borrowing interaction to
non-significance in both waves, though this is not separable from severe collinearity
between coverage and development ($r=0.74$) given only 93 economy clusters. The reversal
itself is therefore observed consistently across both waves, but its most natural specific
explanation is disfavoured by the evidence and its next most natural explanation is not
identifiable with this design --- an empirical puzzle, not a resolved mechanism. It does
not establish that weak coverage causes a negative structural effect, that payment data
are complementary, or that account access, lender processing or supply constraints
explain the result. Future mechanism claims would require lender- or application-level
evidence.

\section*{Limitations}

The evidence is observational repeated cross-sectional evidence, not an individual panel.
The interaction is identified by 93 economy-level historical moderator values, not by
194,558 independent observations of institutional variation. The 2019 coverage measure is
an imperfect historical breadth proxy and does not measure depth, accessibility, legal
usability or lender use. The 2021--2024 payment exposure is qualified-comparable, and 2017
is contextual because its payment construct is not strictly common. The 93-economy frame
is not a random or exhaustive sample of economies with available 2019 coverage data: it
excludes most near-universal-coverage economies because the Findex borrowing module was
not administered there (see Data), compressing the moderator's within-sample range and
narrowing the population to which the boundary-condition language applies. An
inverse-probability-reweighting check (\S5.4) confirms this is not merely a theoretical
concern: 80 of the 132 economies with usable covariates fall outside common support between
the included and excluded groups, so the sample is classified not credibly transportable
with the available covariate overlap to the excluded, mostly high-coverage population,
regardless of the reweighted point estimates' own stability. The same diagnostic also shows the coverage interaction attenuates by 31--39\%
and loses conventional significance once digital-payment activity is additionally allowed
to interact with log GDP per capita and the account-ownership rate. Coverage and log GDP
per capita are themselves highly correlated across the 93 economies ($r=0.74$), so this
attenuation is consistent with either genuine development-level confounding or with
collinearity-driven loss of precision from asking 93 clusters to identify two highly
overlapping economy-level dimensions at once; this design cannot adjudicate between the two,
so the coverage interaction's specificity to credit-information coverage, rather than
economic development more broadly, is not identified here. The 2019
coverage series' Doing Business provenance also carries a disclosed, unresolved
measurement-integrity question for a subset of economies, including two in this sample
(see Data); excluding those two economies (China, Saudi Arabia) and re-standardizing the
moderator over the remaining 91 leaves the primary interaction, the informal-borrowing
placebo, and the macro-control interaction check all qualitatively unchanged (\S3, \S5.4),
so this is a disclosed provenance concern rather than one shown to drive any reported
result. A placebo check on informal borrowing (\S5.4) does not pass: the coverage
interaction is positive, not null, in both waves (2021 $+1.95$ pp, 2024 $+1.75$ pp,
both wild-cluster $p<0.01$), which further weakens any reading of the formal-borrowing
interaction as confirming a credit-information-specific substitution effect. Account
ownership may be a mediator, collider or selection variable. Realised
borrowing does not distinguish demand from supply. The 93-economy, 2021--2024 sample and specification were developed
during exploratory analysis of the data rather than fixed by a prospective, hash-frozen
pre-registration; no such pre-registration exists for this design, and the frozen
\texttt{H-000501} protocol that otherwise governs this project explicitly places
multi-wave/panel estimation of the outcome out of scope pending a new pre-registration
process, which this analysis did not undergo. The reported estimates should therefore be
read as exploratory/observational evidence rather than a confirmatory pre-registered test;
they are independently reproduced at the level of arithmetic (Stage 12) but not
confirmatory in the pre-registration sense. Agreement between AI implementations is
computational reproducibility rather than independent human replication. These
limitations are not resolved by persistence across waves or by robustness checks.

\section*{Data availability, ethics, and conflicts of interest}

\paragraph{Data availability.} The analysis uses the publicly available Global Findex
Database 2021 and 2025 microdata \citep{WorldBankFindexMicrodata2025} and the archived
World Bank Doing Business \emph{Getting Credit} indicators
\citep{WorldBankDoingBusinessArchive}, both obtainable from the World Bank Microdata
Library and Doing Business archive respectively. The frozen analytical dataset (194,558
observations, 93 economies, SHA-256
\texttt{78938AFF518F3F3BC049A18E7B8CA5C1042996392B1ACDD6EAAB56FC555240C1}) and all
estimation code are retained in the project repository and available from the
corresponding author on reasonable request.

\paragraph{Ethics.} All microdata are secondary, de-identified, publicly released survey
data with no individual identifiers; no new human-subjects data were collected for this
study, and no separate ethics/IRB approval was required.

\paragraph{Conflicts of interest.} The author declares no financial or non-financial
conflicts of interest related to this work.

\section*{AI-use disclosure}

AI systems assisted with data engineering, estimation, auditing and typesetting under
author direction. Cross-agent agreement is reported as computational reproducibility, not
independent human replication. The author selected the question, approved the contracts,
reviewed the outputs and remains responsible for all claims and errors. No AI system is an
author.

\section*{Research provenance, exploratory status, and reproducibility}

The 93-economy, 2021--2024 common-frame sample and specification were developed during
exploratory analysis of the underlying Global Findex data, not fixed in advance by a
prospective, hash-frozen pre-registration. No pre-registration record exists for this
design in the project's hypothesis ledger, and the frozen \texttt{H-000501} protocol that
otherwise governs this project explicitly lists multi-wave/panel estimation of the outcome
as out of scope pending a genuinely new pre-registration process, which was not carried
out before this specification was chosen and estimated. This paper's reported associations
are therefore exploratory/observational evidence, not a confirmatory pre-registered test,
and should be read accordingly.

What was frozen, and when, is nonetheless disclosed precisely. The analysis contract and
analytical dataset were hash-recorded after the specification was chosen: the frozen
analytical dataset contains 194,558 observations in 93 economies and has SHA-256
\texttt{78938AFF518F3F3BC049A18E7B8CA5C1042996392B1ACDD6EAAB56FC555240C1}. Stage 5 and
Stage 5.1 result artifacts record the estimator, weights, clustering, bootstrap settings,
specification hierarchy and result checksums, so the arithmetic reported here is fully
traceable to those immutable artifacts; no new estimation was performed for this rewrite.
Tables~\ref{tab:stage8primary} and~\ref{tab:stage8pooled} are generated directly from
those Stage 5/5.1 artifacts by \texttt{code/study\_05\_crosscountry/tables/}
\texttt{build\_stage8\_tables.py} (Stage 15); no coefficient, standard error, $p$-value or
sample count in either table is typed by hand.
An independent forensic replication (Stage 12) rewrote the data-construction and
estimation code from scratch, without reusing any of the code above, and closely
reproduced the primary 2021 and 2024 coefficients, the pooled and triple-interaction
estimates, and all 186 leave-one-economy-out signs. This establishes computational
reproducibility of the numbers reported; it does not establish that the sample or
specification was locked before results were seen, which remains a separate, unresolved
provenance question and is stated plainly above rather than obscured by the reproducibility
result.
